% !PRINTED 509
\weberpage{509}

Um ihn zu beweisen, bezeichnen wir mit $t$ eine in den $\alpha$ und in $\omega$ nicht vorkommende Variable und setzen

\weq{11}{\varphi(t) = t^m + \alpha_1 t^{m-1} + \cdots + \alpha_m.}

Dann ist, wenn $n$ der Grad des Körpers $\Omega$ ist,
\[
\Phi(t) = N[\varphi(t)]
\]
eine ganze Function $mn^{\text{ten}}$ Grades von $t$, deren Coëfficienten ganze rationale Functionale sind. Zugleich ist $\Phi(\omega) = 0$, und folglich $\omega$ eine ganze Zahl.

\gesperrt{Daraus folgt noch, dass der Satz 6. richtig bleibt, wenn die Coëfficienten $A_{k,i}$ nicht ganze Functionale in $R$, sondern in $\Omega$ sind.}

Ist ein ganzes Functional $\omega$ zugleich eine Zahl, so ist es eine ganze Zahl, weil in diesem Falle $N(t-\omega)$ ganze rationale Zahlen zu Coëfficienten hat. Die ganzen Zahlen, die wir im \S.\,133 betrachtet haben, sind demnach als specielle Fälle unter den ganzen Functionalen enthalten.

\S.\,138.

\gesperrt{Theilbarkeit. Associirte Functionale. Einheiten.}

Die ganzen Functionale unterliegen ähnlichen Gesetzen der Theilbarkeit, wie die ganzen rationalen Zahlen. Um sie zu erkennen, stellen wir folgende Definition an die Spitze:

1. \gesperrt{Ein ganzes Functional $\alpha$ heisst durch ein anderes, von Null verschiedenes ganzes Functional $\beta$ theilbar, wenn der Quotient $\alpha:\beta=\gamma$ ein ganzes Functional ist.}

Es ist dann $\alpha=\beta\gamma$ und $\beta$ und $\gamma$ heissen \gesperrt{Theiler} von $\alpha$, und man sagt auch, \gesperrt{$\beta$ und $\gamma$ gehen in $\alpha$ auf}. Die Zahl 0 ist durch \gesperrt{jedes ganze} Functional theilbar.

Aus dieser Definition ergeben sich ohne Schwierigkeit die folgenden fundamentalen Sätze über Theilbarkeit:

2. \gesperrt{Sind $\alpha$ und $\alpha_1$ theilbar durch $\beta$, so ist auch $\alpha\pm\alpha_1$ theilbar durch $\beta$.}

Denn ist $\alpha=\beta\gamma$, $\alpha_1=\beta\gamma_1$, so ist $\alpha\pm\alpha_1=\beta(\gamma\pm\gamma_1)$, nnd wenn $\gamma$ und $\gamma_1$ ganz sind, so ist auch $\gamma\pm\gamma_1$ ganz. % [sic] printed ``nnd''
